\documentclass[11pt,letterpaper]{article}

\usepackage[margin=1in]{geometry}
\usepackage[T1]{fontenc}
\usepackage{lmodern}
\usepackage{amsmath,amssymb}
\usepackage{booktabs}
\usepackage{microtype}
\usepackage[hidelinks]{hyperref}

\setlength{\parindent}{0pt}
\setlength{\parskip}{6pt}
\newcommand{\sinc}{\operatorname{sinc}}

\title{Fourier Transform:\\
Phase, Sampling, and Time--Frequency Localization}
\author{Signal Processing}
\date{}

\begin{document}
\maketitle

The Fourier transform describes a signal in terms of its frequency
components. This handout distinguishes three related ideas:
the amplitude and phase of a component, the sampling of the time and
frequency axes, and the localization of a signal in these domains.

\section{Continuous Fourier transform}

We use frequency $f$ in hertz and angular frequency
$\omega=2\pi f$. Our transform convention is
\begin{equation}
\boxed{
X(f)=\int_{-\infty}^{\infty}x(t)e^{-i2\pi ft}\,dt,
\qquad
x(t)=\int_{-\infty}^{\infty}X(f)e^{i2\pi ft}\,df.
}
\end{equation}

Euler's identity,
\[
e^{i\theta}=\cos\theta+i\sin\theta,
\]
connects the complex exponential representation with cosine and sine.

For a real signal,
\[
X(-f)=X(f)^*,
\]
where the asterisk denotes complex conjugation. Thus, positive and
negative frequencies are related rather than independent.

\section{Cosine--sine coefficients and amplitude--phase}

At one nonzero frequency, a real sinusoidal component can be written as
\begin{equation}
x_\omega(t)=a\cos(\omega t)+b\sin(\omega t).
\end{equation}
The coefficients $a$ and $b$ are two independent real numbers.

Using
\[
\cos(\omega t-\phi)
=\cos\phi\cos(\omega t)+\sin\phi\sin(\omega t),
\]
the same component becomes
\begin{equation}
\boxed{x_\omega(t)=A\cos(\omega t-\phi)}
\end{equation}
with
\begin{equation}
\boxed{
a=A\cos\phi,\qquad b=A\sin\phi,\qquad
A=\sqrt{a^2+b^2},\qquad
\phi=\operatorname{atan2}(b,a).
}
\end{equation}
If $A=0$, phase is undefined.

These are Cartesian and polar descriptions of the same coefficient
vector: $(a,b)$ gives its coordinates, while $(A,\phi)$ gives its
length and direction.

\subsection*{Why are cosine and sine $90^\circ$ apart?}

The fixed basis functions satisfy
\[
\sin(\omega t)=\cos\left(\omega t-\frac{\pi}{2}\right).
\]
They are therefore \emph{in quadrature}: their relative phase is
$90^\circ$.

The phase $\phi$ describes the resulting sinusoid relative to the
time origin. It does not change the relative phase of the fixed
basis functions.

One could introduce separate shifts,
\[
a\cos(\omega t+\alpha)+b\sin(\omega t+\beta),
\]
but expanding gives
\[
(a\cos\alpha+b\sin\beta)\cos(\omega t)
+(-a\sin\alpha+b\cos\beta)\sin(\omega t).
\]
Thus, separate shifts produce another cosine--sine combination;
the four parameters are redundant.

\subsection*{Connection with a complex Fourier coefficient}

With the synthesis exponential $e^{i\omega t}$,
\[
x_\omega(t)=\operatorname{Re}\{C e^{i\omega t}\},
\qquad C=a-ib.
\]
Consequently,
\[
|C|=A,\qquad \arg C=-\phi.
\]
The minus sign follows from our choice
$x_\omega(t)=A\cos(\omega t-\phi)$.

For a Fourier series written using both positive and negative
frequencies, the corresponding coefficients are
\[
c_+=\frac{a-ib}{2},\qquad c_-=\frac{a+ib}{2},
\]
because
\[
x_\omega(t)=c_+e^{i\omega t}+c_-e^{-i\omega t}.
\]

\section{Sampling, frequency range, and frequency spacing}

Use different symbols for sampling intervals and signal widths.

\begin{center}
\begin{tabular}{lll}
\toprule
Symbol & Meaning & Units\\
\midrule
$\delta t$ & Time sampling interval & s\\
$\delta f$ & DFT frequency-bin spacing & Hz\\
$T$ & DFT record duration & s\\
$f_N$ & Nyquist frequency & Hz\\
$\sigma_t$ & Energy-based temporal spread & s\\
$\sigma_f$ & Energy-based spectral spread & Hz\\
\bottomrule
\end{tabular}
\end{center}

For $N$ samples taken at $t_n=n\delta t$, $n=0,\ldots,N-1$,
the DFT period is
\[
T=N\delta t.
\]
The elapsed time between the first and last samples is
$(N-1)\delta t$; $N\delta t$ is the duration used in the DFT
spacing relation.

The DFT frequency spacing is
\begin{equation}
\boxed{
\delta f=\frac{1}{N\delta t}=\frac{1}{T},
\qquad N\delta t\,\delta f=1.
}
\end{equation}

The Nyquist frequency is
\begin{equation}
\boxed{f_N=\frac{1}{2\delta t}.}
\end{equation}
To avoid aliasing, the sampled signal must have its spectral content
restricted below this limit, normally through appropriate filtering.

These equations express different controls:
\begin{itemize}
\item Smaller $\delta t$ increases the available frequency range.
\item Longer observation duration $T$ decreases the frequency-bin
spacing.
\end{itemize}

The full DFT frequency range has width
\[
N\delta f=\frac{1}{\delta t}=2f_N.
\]
When expressed using signed frequencies, this is a frequency interval
of width $2f_N$, conventionally chosen around zero.

\subsection*{Example}

With
\[
\delta t=0.001\ {\rm s},\qquad N=1000,
\]
we obtain
\[
T=1\ {\rm s},\qquad
\delta f=1\ {\rm Hz},\qquad f_N=500\ {\rm Hz}.
\]
Recording for $2$ s at the same sampling interval gives
$\delta f=0.5$ Hz, while $f_N$ remains $500$ Hz.

\subsection*{Frequency spacing versus resolving power}

Frequency-bin spacing describes the DFT grid. The ability to
distinguish nearby spectral components also depends on observation
duration, window shape, noise, and the estimation method.

For a rectangular observation window, the spectral main-lobe width
scales as $1/T$. Increasing the duration of actual observations can
therefore improve frequency discrimination.

Appending zeros decreases the DFT bin spacing but adds no observations.
It samples the existing spectrum more densely; it does not narrow
the spectral main lobe or improve the resolving power of the data.

\section{Fourier transform of a boxcar}

Consider a unit-height boxcar of duration $L$:
\[
x(t)=
\begin{cases}
1,& |t|\leq L/2,\\
0,& |t|>L/2.
\end{cases}
\]
We use $L$ for the pulse duration to distinguish it from the DFT
record duration $T$.

Its Fourier transform is
\begin{align}
X(f)
&=\int_{-L/2}^{L/2}e^{-i2\pi ft}\,dt\\
&=\left[\frac{e^{-i2\pi ft}}{-i2\pi f}\right]_{-L/2}^{L/2}\\
&=\frac{e^{-i\pi fL}-e^{i\pi fL}}{-i2\pi f}\\
&=\frac{\sin(\pi fL)}{\pi f}.
\end{align}
Define the normalized sinc function by
\[
\sinc(u)=\frac{\sin(\pi u)}{\pi u},\qquad \sinc(0)=1.
\]
Then
\begin{equation}
\boxed{
\operatorname{rect}(t/L)
\quad\longleftrightarrow\quad
L\,\sinc(fL).
}
\end{equation}
At $f=0$, the limit gives $X(0)=L$, equal to the area of the boxcar.

The first spectral zeros occur at $f=\pm1/L$, so the main-lobe
width between these zeros is
\begin{equation}
\boxed{W_{f,\mathrm{main}}=\frac{2}{L}.}
\end{equation}
A longer boxcar produces a narrower main lobe. Its abrupt edges
produce sidelobes whose amplitude decays as $1/|f|$.

Multiplication by a boxcar observation window convolves a signal's
spectrum with this sinc function. This explains spectral broadening
and leakage associated with finite rectangular records.

\section{Fourier transform of a Gaussian}

Let
\[
x(t)=e^{-t^2/(2s^2)},\qquad s>0.
\]
Here $s$ is the width parameter of the \emph{amplitude} Gaussian;
it is not yet the energy-based standard deviation $\sigma_t$.

Its transform is
\[
X(f)=\int_{-\infty}^{\infty}
e^{-t^2/(2s^2)}e^{-i2\pi ft}\,dt.
\]

We can derive the result by differentiating with respect to $f$.
Since
\[
x'(t)=-\frac{t}{s^2}x(t),
\]
we obtain
\begin{align}
X'(f)
&=-i2\pi\int_{-\infty}^{\infty}
t\,x(t)e^{-i2\pi ft}\,dt\\
&=i2\pi s^2\int_{-\infty}^{\infty}
x'(t)e^{-i2\pi ft}\,dt.
\end{align}
Integration by parts gives
\[
\int_{-\infty}^{\infty}x'(t)e^{-i2\pi ft}\,dt
=i2\pi fX(f),
\]
because the boundary term vanishes. Therefore,
\[
X'(f)=-4\pi^2s^2fX(f).
\]
Solving this differential equation,
\[
X(f)=X(0)e^{-2\pi^2s^2f^2}.
\]
The Gaussian integral gives
\[
X(0)=\int_{-\infty}^{\infty}e^{-t^2/(2s^2)}\,dt
=\sqrt{2\pi}\,s.
\]
Thus,
\begin{equation}
\boxed{
e^{-t^2/(2s^2)}
\quad\longleftrightarrow\quad
\sqrt{2\pi}\,s\,e^{-2\pi^2s^2f^2}.
}
\end{equation}

The transform of a Gaussian is another Gaussian. Writing the
frequency amplitude as $C e^{-f^2/(2s_f^2)}$ gives
\[
s_f=\frac{1}{2\pi s},
\qquad
\boxed{s\,s_f=\frac{1}{2\pi}.}
\]
These are amplitude-profile width parameters.

\section{Time--frequency uncertainty}

The uncertainty principle concerns the spread of the
\emph{signal's energy}, rather than the spacing of sampled axes.

For a nonzero finite-energy signal, define
\[
E=\int_{-\infty}^{\infty}|x(t)|^2\,dt
=\int_{-\infty}^{\infty}|X(f)|^2\,df,
\]
where the equality follows from Parseval's theorem.

The energy centroids are
\[
\bar t=\frac{1}{E}\int t|x(t)|^2\,dt,
\qquad
\bar f=\frac{1}{E}\int f|X(f)|^2\,df.
\]
The corresponding variances are
\begin{equation}
\sigma_t^2=\frac{1}{E}\int(t-\bar t)^2|x(t)|^2\,dt,
\qquad
\sigma_f^2=\frac{1}{E}\int(f-\bar f)^2|X(f)|^2\,df.
\end{equation}
All integrals above extend over the real line.

When these variances are finite, the Fourier uncertainty relation is
\begin{equation}
\boxed{\sigma_t\sigma_f\geq\frac{1}{4\pi}.}
\end{equation}
In angular frequency, $\sigma_\omega=2\pi\sigma_f$, so
\[
\boxed{\sigma_t\sigma_\omega\geq\frac12.}
\]

A signal cannot have arbitrarily small energy spreads in both time
and frequency.

\subsection*{The Gaussian reaches equality}

For $x(t)=e^{-t^2/(2s^2)}$,
\[
|x(t)|^2=e^{-t^2/s^2},
\]
so
\[
\sigma_t=\frac{s}{\sqrt2}.
\]
Similarly,
\[
|X(f)|^2=2\pi s^2e^{-4\pi^2s^2f^2},
\]
giving
\[
\sigma_f=\frac{1}{2\sqrt2\,\pi s}.
\]
Therefore,
\begin{equation}
\boxed{
\sigma_t\sigma_f
=\frac{s}{\sqrt2}\,
\frac{1}{2\sqrt2\,\pi s}
=\frac{1}{4\pi}.
}
\end{equation}
This differs from the amplitude-width product $s\,s_f=1/(2\pi)$
because squaring a Gaussian makes its profile narrower.

\subsection*{The boxcar has infinite spectral variance}

For the boxcar,
\[
\sigma_t=\frac{L}{\sqrt{12}}.
\]
However,
\[
f^2|X(f)|^2=\frac{\sin^2(\pi fL)}{\pi^2},
\]
whose integral over all frequencies diverges. Therefore,
\[
\sigma_f=\infty.
\]
The sinc main lobe has finite width, but its energy-based frequency
variance is infinite. Main-lobe width and standard deviation are
different measures of spectral width.

\section{Key distinctions}

\begin{center}
\begin{tabular}{lp{0.62\textwidth}}
\toprule
Relation & Interpretation\\
\midrule
$f_N=1/(2\delta t)$ &
Time sampling controls the available frequency range.\\[4pt]
$\delta f=1/T$ &
DFT record duration controls frequency-bin spacing.\\[4pt]
$\sigma_t\sigma_f\geq1/(4\pi)$ &
Signal energy cannot be arbitrarily localized in both domains.\\
\bottomrule
\end{tabular}
\end{center}

A short pulse can be sampled very finely and placed in a long,
zero-padded record. Its DFT bins can then be closely spaced, while
its intrinsic spectrum remains broad.

\begin{quote}
Cosine and sine are fixed basis functions in quadrature.
Amplitude and phase describe their resulting sinusoid.
Sampling intervals describe the axes; energy spreads describe
the signal.
\end{quote}

\end{document}